\section{Numerical studies}\label{sec:numerics}
All computations use the USEM-MATLAB code base (GNU Octave~8.4, IEEE double precision), in
which the scaled Lagrange basis is the default of the metric space $Q_p(\hat x,\hat y)$
(\PM) and the hybrid space is available as a second trial space (\PH); \PP\ is the
conventional spectral element on the same nodes. GLL nodes are used throughout. The scripts
and data are listed in \cref{sec:reproduction}.

\subsection{Test elements and meshes}\label{sec:elements}
\emph{Single elements:} a parallelogram (affine), a trapezoid with straight edges (bilinear
map), an element with two curved edges (``curved'') and one with angular and curved-edge
distortion (``general''), $p=1,\dots,16$. \emph{Mesh level:} a $3\times2$ mesh on
$[0,3]\times[0,2]$ with curved-edge distortion $0.15$, $0.20$ and combined angular and
curved-edge distortion $0.30$ (minimum $\det\bm J=0.16$, $0.09$, $0.014$). \emph{Patch tests:}
the square panel of \citet[Secs.~5.5--5.7]{Rajendran2010technique}, $2\times2$ elements whose
internal mid-side nodes are rotated by $\theta=90^\circ$--$270^\circ$ (every element has
$\det\bm J\le0$ somewhere), with interior nodes from the bilinear map of the corners.
\emph{Convergence:} the thick cylinder of \citet[Sec.~5.9]{Rajendran2010technique} (quarter
annulus $1\le r\le3$, plane stress, $E=1000$, $\nu=0.25$, internal pressure $1000$),
$2N\times2N$ elements, $N=1,2,4,8$, generated with gmsh~\cite{Geuzaine2009Gmsh}; in each
$2\times2$ super-element the internal mid-side nodes are rotated by
$\alpha\in\{0^\circ,15^\circ,30^\circ,45^\circ,60^\circ,90^\circ\}$, so that the
\emph{relative} distortion is the same on all meshes and the element maps are not asymptotically
affine (\cref{fig:meshes}). Folded elements occur for $\alpha=45^\circ$ ($N\ge2$) and
$\alpha\ge60^\circ$. The arcs are represented exactly (Gordon--Hall blending~\cite{Gordon1973Transfinite}).
Three fields are computed: the Lam\'e solution, a manufactured polynomial of total degree~5, and
the analytic field $u_x=\tfrac1{10}(\sin3X\cos2Y+e^{XY})$, $u_y=\tfrac1{10}(\cos2X\sin3Y+Ye^{-X})$,
$X=x/3$, $Y=y/3$, with body forces and tractions derived from the exact fields.
Errors: relative energy norm $\eE$ (strains of the respective trial functions, signed
$\det\bm J$) and relative maximum nodal error $e_n$. Cases in which the code base finds the
interpolation matrix singular to working precision ($\mathrm{rcond}<\varepsilon$) are marked
``rej.''.

\begin{figure}[t]
  \centering
  \begin{tikzpicture}
    \begin{groupplot}[group style={group size=6 by 1, horizontal sep=1mm}, meshplot,
        xmin=-0.05, xmax=2.05, ymin=-1.05, ymax=1.05, width=3.0cm, height=3.0cm]
      \nextgroupplot[title={$\theta=90^\circ$}]
      \addplot[black, line width=0.3pt] table {\data/panel_mesh_N1_t090.dat};
      \nextgroupplot[title={$\theta=135^\circ$}]
      \addplot[black, line width=0.3pt] table {\data/panel_mesh_N1_t135.dat};
      \nextgroupplot[title={$\theta=175^\circ$}]
      \addplot[black, line width=0.3pt] table {\data/panel_mesh_N1_t175.dat};
      \nextgroupplot[title={$\theta=225^\circ$}]
      \addplot[black, line width=0.3pt] table {\data/panel_mesh_N1_t225.dat};
      \nextgroupplot[title={$\theta=270^\circ$}]
      \addplot[black, line width=0.3pt] table {\data/panel_mesh_N1_t270.dat};
      \nextgroupplot[title={$\theta=0^\circ$}]
      \addplot[black, line width=0.3pt] table {\data/panel_mesh_N1_t000.dat};
    \end{groupplot}
  \end{tikzpicture}\\[1mm]
  \begin{tikzpicture}
    \begin{groupplot}[group style={group size=6 by 1, horizontal sep=1mm}, meshplot,
        xmin=-0.05, xmax=3.05, ymin=-0.05, ymax=3.05, width=3.0cm, height=3.0cm]
      \nextgroupplot[title={$\alpha=0^\circ$}]
      \addplot[black, line width=0.3pt] table {\data/cyl_mesh_N2_a000.dat};
      \nextgroupplot[title={$\alpha=15^\circ$}]
      \addplot[black, line width=0.3pt] table {\data/cyl_mesh_N2_a015.dat};
      \nextgroupplot[title={$\alpha=30^\circ$}]
      \addplot[black, line width=0.3pt] table {\data/cyl_mesh_N2_a030.dat};
      \nextgroupplot[title={$\alpha=45^\circ$}]
      \addplot[black, line width=0.3pt] table {\data/cyl_mesh_N2_a045.dat};
      \nextgroupplot[title={$\alpha=60^\circ$}]
      \addplot[black, line width=0.3pt] table {\data/cyl_mesh_N2_a060.dat};
      \nextgroupplot[title={$\alpha=90^\circ$}]
      \addplot[black, line width=0.3pt] table {\data/cyl_mesh_N2_a090.dat};
    \end{groupplot}
  \end{tikzpicture}
  \caption{Top: panel meshes of the Jacobian-unlimited patch test (internal mid-side nodes
    rotated by $\theta$). Bottom: thick-cylinder meshes for $N=2$ ($4\times4$ elements); the finer
    meshes repeat the pattern on $N^2$ super-elements.}
  \label{fig:meshes}
\end{figure}

\subsection{Element level: conditioning and Lebesgue constants}\label{sec:num-elements}
\Cref{fig:basis} shows $\kappa_2(\bm V)$ and the error of the computed shape functions for
worst-case (unit) nodal data, $e_\delta=\max_{k,l}|M_k(\bx_l)-\delta_{kl}|$, for the space
$Q_p(\hat x,\hat y)$ in four bases. With monomials $\kappa_2$ grows by a factor $\approx5.7$
per degree already on the parallelogram ($3.5\cdot10^{11}$ at $p=16$), close to
$(1+\sqrt2)^2$, i.e.\ $(1+\sqrt2)^p$ for each 1D factor~\cite{Gautschi1988Lower}; the scaled
Lagrange basis gives $\kappa_2=1$ for every $p$ (\cref{lem:affine}), and the element matrix is
symmetric to round-off, whereas its asymmetry grows to $2\cdot10^{-6}$ with monomials
(\cref{tab:affine}). On the distorted elements the three alternatives lie within a factor~3 of
one another and are better than monomials by factors $4.6\cdot10^3$--$2.9\cdot10^4$ at $p=16$; $e_\delta$
follows $\kappa_2u$ and decreases by factors $3.6\cdot10^2$--$8.2\cdot10^3$ at $p=16$ (e.g.\
trapezoid: $3.0\cdot10^{-2}\to7.1\cdot10^{-6}$). For $p\le6$ all bases give $e_\delta<5\cdot10^{-13}$.
Without the bounding-box scaling~\eqref{eq:bbox} the orthogonal and Lagrange bases lose more than three
orders of magnitude on the trapezoid ($\kappa_2=2.2\cdot10^{15}$ instead of $8.6\cdot10^{11}$ at
$p=16$).

\begin{figure}[t]
  \centering
  \begin{tikzpicture}
    \begin{groupplot}[group style={group size=3 by 2, horizontal sep=13mm, vertical sep=16mm},
        width=5.2cm, height=4.4cm, ymode=log, xmin=1, xmax=16, xlabel={$p$},
        every axis plot/.append style={line width=0.7pt, mark size=1.1pt},
        unbounded coords=discard]
      \nextgroupplot[title={trapezoid}, ylabel={$\kappa_2(\bm V)$}, ymin=0.5, ymax=1e16,
        legend to name=leg:basis, legend columns=4, legend style={font=\footnotesize}]
      \addplot[cB1, mark=*] table[x=p, y=kappa_mon]{\data/metric-basis-trapezoid.dat};
      \addplot[cB2, mark=square*] table[x=p, y=kappa_leg]{\data/metric-basis-trapezoid.dat};
      \addplot[cB3, mark=triangle*] table[x=p, y=kappa_che]{\data/metric-basis-trapezoid.dat};
      \addplot[cB4, mark=diamond*] table[x=p, y=kappa_lag]{\data/metric-basis-trapezoid.dat};
      \legend{monomial, Legendre (scaled), Chebyshev (scaled), Lagrange (scaled)}
      \nextgroupplot[title={curved}, ymin=0.5, ymax=1e16]
      \addplot[cB1, mark=*] table[x=p, y=kappa_mon]{\data/metric-basis-curved.dat};
      \addplot[cB2, mark=square*] table[x=p, y=kappa_leg]{\data/metric-basis-curved.dat};
      \addplot[cB3, mark=triangle*] table[x=p, y=kappa_che]{\data/metric-basis-curved.dat};
      \addplot[cB4, mark=diamond*] table[x=p, y=kappa_lag]{\data/metric-basis-curved.dat};
      \nextgroupplot[title={general}, ymin=0.5, ymax=1e16]
      \addplot[cB1, mark=*] table[x=p, y=kappa_mon]{\data/metric-basis-general.dat};
      \addplot[cB2, mark=square*] table[x=p, y=kappa_leg]{\data/metric-basis-general.dat};
      \addplot[cB3, mark=triangle*] table[x=p, y=kappa_che]{\data/metric-basis-general.dat};
      \addplot[cB4, mark=diamond*] table[x=p, y=kappa_lag]{\data/metric-basis-general.dat};
      \nextgroupplot[ylabel={$e_\delta$}, ymin=1e-17, ymax=1]
      \addplot[cB1, mark=*] table[x=p, y=eKron_mon]{\data/metric-basis-trapezoid.dat};
      \addplot[cB2, mark=square*] table[x=p, y=eKron_leg]{\data/metric-basis-trapezoid.dat};
      \addplot[cB3, mark=triangle*] table[x=p, y=eKron_che]{\data/metric-basis-trapezoid.dat};
      \addplot[cB4, mark=diamond*] table[x=p, y=eKron_lag]{\data/metric-basis-trapezoid.dat};
      \nextgroupplot[ymin=1e-17, ymax=1]
      \addplot[cB1, mark=*] table[x=p, y=eKron_mon]{\data/metric-basis-curved.dat};
      \addplot[cB2, mark=square*] table[x=p, y=eKron_leg]{\data/metric-basis-curved.dat};
      \addplot[cB3, mark=triangle*] table[x=p, y=eKron_che]{\data/metric-basis-curved.dat};
      \addplot[cB4, mark=diamond*] table[x=p, y=eKron_lag]{\data/metric-basis-curved.dat};
      \nextgroupplot[ymin=1e-17, ymax=1]
      \addplot[cB1, mark=*] table[x=p, y=eKron_mon]{\data/metric-basis-general.dat};
      \addplot[cB2, mark=square*] table[x=p, y=eKron_leg]{\data/metric-basis-general.dat};
      \addplot[cB3, mark=triangle*] table[x=p, y=eKron_che]{\data/metric-basis-general.dat};
      \addplot[cB4, mark=diamond*] table[x=p, y=eKron_lag]{\data/metric-basis-general.dat};
    \end{groupplot}
  \end{tikzpicture}\\[1mm]
  \ref{leg:basis}
  \caption{Space $Q_p(\hat x,\hat y)$ on the distorted test elements: condition number of the
    interpolation matrix (top) and error of the Kronecker property of the computed shape
    functions (bottom) for four bases.}
  \label{fig:basis}
\end{figure}

\begin{table}[t]
  \centering\small
  \caption{Parallelogram (affine element), $Q_p(\hat x,\hat y)$: asymmetry
    $\|\bm K^e-\bm K^{e\T}\|_F/\|\bm K^e\|_F$ of the element matrix (zero in exact arithmetic) and
    condition number of $\bm V$.}
  \label{tab:affine}
  \begin{tabular}{@{}lcccccc@{}}
\toprule
$p$ & 4 & 6 & 8 & 10 & 12 & 16\\
\midrule
monomial & $2.2\cdot10^{-15}$ & $8.4\cdot10^{-14}$ & $4.2\cdot10^{-12}$ & $1.3\cdot10^{-10}$ & $2.9\cdot10^{-9}$ & $2.0\cdot10^{-6}$\\
Lagrange (scaled) & $1.6\cdot10^{-15}$ & $2.0\cdot10^{-15}$ & $4.6\cdot10^{-15}$ & $1.5\cdot10^{-14}$ & $1.5\cdot10^{-14}$ & $3.4\cdot10^{-14}$\\
\midrule
$\kappa_2(\bm V)$, monomial & $2.9\cdot10^{2}$ & $9.2\cdot10^{3}$ & $2.9\cdot10^{5}$ & $9.6\cdot10^{6}$ & $3.2\cdot10^{8}$ & $3.5\cdot10^{11}$\\
$\kappa_2(\bm V)$, Lagrange (scaled) & $1.0$ & $1.0$ & $1.0$ & $1.0$ & $1.0$ & $1.0$\\
\bottomrule
\end{tabular}

\end{table}

As \cref{prop:invariance} predicts, the reproduction of polynomials and the Lebesgue constant
are the same in all bases. \Cref{fig:lebesgue} compares $\Lambda$ of $Q_p(\hat x,\hat y)$ and of
the hybrid space $T_p$ on the three distorted elements. For $Q_p(\hat x,\hat y)$, $\Lambda$ grows
exponentially, to $4.0\cdot10^7$ (trapezoid), $5.5\cdot10^4$ (curved) and $6.7\cdot10^6$ (general) at
$p=16$. For $T_p$ it equals $\Lambda_{\mathrm{GLL}}$ on the trapezoid (bilinear map,
\cref{prop:structure}) and stays within $2\,\%$ (curved) and $40\,\%$ (general) of it ($6.1$ and
$8.5$ at $p=16$).

\begin{figure}[t]
  \centering
  \begin{tikzpicture}
    \begin{semilogyaxis}[width=8.5cm, height=5.6cm, xlabel={$p$}, ylabel={Lebesgue constant},
        xmin=2, xmax=16, ymax=1e9, legend columns=2,
        legend style={font=\scriptsize, at={(0.02,0.98)}, anchor=north west},
        every axis plot/.append style={line width=0.7pt, mark size=1.3pt}]
      \addplot[cPM, mark=*] table[x=p, y=LebQ_trapezoid]{\data/metric-basis-hybrid.dat};
      \addplot[cPH, mark=*, dashed] table[x=p, y=LebH_trapezoid]{\data/metric-basis-hybrid.dat};
      \addplot[cPM, mark=square*] table[x=p, y=LebQ_curved]{\data/metric-basis-hybrid.dat};
      \addplot[cPH, mark=square*, dashed] table[x=p, y=LebH_curved]{\data/metric-basis-hybrid.dat};
      \addplot[cPM, mark=triangle*] table[x=p, y=LebQ_general]{\data/metric-basis-hybrid.dat};
      \addplot[cPH, mark=triangle*, dashed] table[x=p, y=LebH_general]{\data/metric-basis-hybrid.dat};
      \addplot[black, dotted, line width=1pt] table[x=p, y=LebN]{\data/metric-basis-parallelogram.dat};
      \legend{$Q_p$ trapezoid, $T_p$ trapezoid, $Q_p$ curved, $T_p$ curved, $Q_p$ general,
        $T_p$ general, GLL}
    \end{semilogyaxis}
  \end{tikzpicture}
  \caption{Lebesgue constant of the metric interpolation in $Q_p(\hat x,\hat y)$ and in the hybrid
    space $T_p$ on the distorted test elements; GLL: parametric interpolation (any element).}
  \label{fig:lebesgue}
\end{figure}

\subsection{Mesh level: smooth field and eigenvalues}\label{sec:num-mesh}
On the distorted $3\times2$ meshes (\cref{fig:mesh}) the patch test of degree $p$ (harmonic
polynomial field) fails for the conventional element ($10^{-7}$--$10^{-2}$) and is passed by both
unsymmetric elements, with errors growing to $10^{-13}$--$10^{-12}$ at $p=10$ for $Q_p(\hat x,\hat
y)$ and $\le3.1\cdot10^{-15}$ for $T_p$. For the field $\bm u=(e^x\cos y,-e^x\sin y)$ all three elements
converge exponentially; for $p\le8$, $Q_p(\hat x,\hat y)$ is in most cases the more accurate of the
two unsymmetric elements (by factors up to $17$; exceptions: curved $0.20$, $p=6,8$; combined, $p=7$),
at $p=9,10$ $T_p$ is more accurate or within a factor $1.5$, and $T_p$ is $3.6$--$150$ times more accurate
than \PP\ for $p=5,\dots,10$ (curved $0.20$). The decisive difference
is in the spectrum of the generalised eigenvalue problem $\bm K\bm\phi=\lambda\bm M\bm\phi$ (consistent
mass, homogeneous Dirichlet conditions; reference: conventional element on the undistorted mesh,
$p=12$). With $Q_p(\hat x,\hat y)$, from $p=5$--$7$ on $2$--$25$ eigenvalues have negative real part,
the maximum of $|\operatorname{Im}\lambda|/|\lambda|$ approaches $1$, and the error of the eight
lowest eigenvalues stagnates at $10^{-5}$--$10^{-3}$ or reaches $O(1)$ (curved $0.20$, $p=6$: $1.9$).
With $T_p$ no eigenvalue has negative real part, $|\operatorname{Im}\lambda|/|\lambda|\le1.8\cdot10^{-2}$
over the whole spectrum, and the errors of the eight lowest eigenvalues do not exceed those of the
conventional element by more than $6\,\%$ and are smaller in most cases ($4.5\cdot10^{-7}$ for both at
$p=10$). An eigenvalue with negative real part
is an exponentially growing mode for every time integrator; for transient and modal analyses at
$p\gtrsim5$ on curved meshes this disqualifies $Q_p(\hat x,\hat y)$.

\begin{figure}[t]
  \centering
  \begin{tikzpicture}
    \begin{groupplot}[group style={group size=3 by 2, horizontal sep=15mm, vertical sep=17mm},
        width=5.2cm, height=4.4cm, ymode=log, xmin=1, xmax=10, xlabel={$p$},
        every axis plot/.append style={mark size=1.2pt}, unbounded coords=discard]
      \nextgroupplot[title={curved $0.20$}, ylabel={$e_n$, $e^x\cos y$},
        legend to name=leg:mesh, legend columns=3, legend style={font=\footnotesize}]
      \addplot[sPP] table[x=p, y=smooth_cla]{\data/metric-hybrid-curved20.dat};
      \addplot[sPM] table[x=p, y=smooth_skew]{\data/metric-hybrid-curved20.dat};
      \addplot[sPH] table[x=p, y=smooth_hyb]{\data/metric-hybrid-curved20.dat};
      \legend{parametric (PP), {$Q_p$, skew coordinates (PM)}, {$T_p$, hybrid (PH)}}
      \nextgroupplot[title={curved $0.20$}, ylabel={$e_\lambda$ (8 lowest)}]
      \addplot[sPP] table[x=p, y=eig_cla]{\data/metric-hybrid-curved20.dat};
      \addplot[sPM] table[x=p, y=eig_skew]{\data/metric-hybrid-curved20.dat};
      \addplot[sPH] table[x=p, y=eig_hyb]{\data/metric-hybrid-curved20.dat};
      \nextgroupplot[title={curved $0.20$}, ylabel={$\max|\operatorname{Im}\lambda|/|\lambda|$}, ymin=1e-4, ymax=2]
      \addplot[sPM] table[x=p, y=imag_skew]{\data/metric-hybrid-curved20.dat};
      \addplot[sPH] table[x=p, y=imag_hyb]{\data/metric-hybrid-curved20.dat};
      \nextgroupplot[title={angular + curved $0.30$}, ylabel={$e_n$, $e^x\cos y$}]
      \addplot[sPP] table[x=p, y=smooth_cla]{\data/metric-hybrid-combined30.dat};
      \addplot[sPM] table[x=p, y=smooth_skew]{\data/metric-hybrid-combined30.dat};
      \addplot[sPH] table[x=p, y=smooth_hyb]{\data/metric-hybrid-combined30.dat};
      \nextgroupplot[title={angular + curved $0.30$}, ylabel={$e_\lambda$ (8 lowest)}]
      \addplot[sPP] table[x=p, y=eig_cla]{\data/metric-hybrid-combined30.dat};
      \addplot[sPM] table[x=p, y=eig_skew]{\data/metric-hybrid-combined30.dat};
      \addplot[sPH] table[x=p, y=eig_hyb]{\data/metric-hybrid-combined30.dat};
      \nextgroupplot[title={angular + curved $0.30$}, ylabel={$\max|\operatorname{Im}\lambda|/|\lambda|$}, ymin=1e-4, ymax=2]
      \addplot[sPM] table[x=p, y=imag_skew]{\data/metric-hybrid-combined30.dat};
      \addplot[sPH] table[x=p, y=imag_hyb]{\data/metric-hybrid-combined30.dat};
    \end{groupplot}
  \end{tikzpicture}\\[1mm]
  \ref{leg:mesh}
  \caption{$3\times2$ meshes with curved-edge distortion: nodal error for the field
    $(e^x\cos y,-e^x\sin y)$ (left), error of the real parts of the eight eigenvalues of smallest
    modulus (centre), and largest relative imaginary part over the whole spectrum (right; zero for
    the symmetric conventional element).}
  \label{fig:mesh}
\end{figure}

\subsection{Jacobian-unlimited patch test}\label{sec:num-panel}
On the panel meshes the shear load case of \citet{Rajendran2010technique} has a cubic exact
solution, contained in both unsymmetric spaces for $p\ge3$. \Cref{tab:panel} gives the maximum
nodal error over the five distorted meshes (\cref{fig:meshes}). Both unsymmetric elements
reproduce the field on all meshes with folded and overlapping elements; the round-off level of
\PM\ reaches $1.1\cdot10^{-5}$ at $p=8$, that of \PH\ $6.7\cdot10^{-8}$. The conventional element has
errors of $13$--$140\,\%$. For $p=2$ (9-node element, corner-centroid interior node, $5\times5$ Gauss
points) \PH\ passes the constant and linear strain patch tests on all meshes and converges with
the energy-norm rates $2.02$ and $2.21$ on the folded mesh sequences B and C of the paper (PM-QUAD9$^\ast$:
$2.00$, $1.98$; conventional element: $0.70$, $2.50$, the latter from a negative displacement on
the coarsest mesh).

\begin{table}[t]
  \centering\small
  \caption{Jacobian-unlimited patch test (cubic field), maximum relative nodal error over the
    panel meshes with $\theta=90^\circ,\dots,270^\circ$.}
  \label{tab:panel}
  \begin{tabular}{@{}lcccccc@{}}
\toprule
 & $p=3$ & $p=4$ & $p=5$ & $p=6$ & $p=7$ & $p=8$\\
\midrule
parametric (\PP) & $1.6\cdot10^{-1}$ & $4.9\cdot10^{-1}$ & $1.5\cdot10^{-1}$ & $2.2\cdot10^{-1}$ & $1.3\cdot10^{-1}$ & $1.4$\\
$Q_p(\hat x,\hat y)$ (\PM) & $1.3\cdot10^{-13}$ & $3.3\cdot10^{-11}$ & $9.8\cdot10^{-11}$ & $6.5\cdot10^{-11}$ & $1.1\cdot10^{-9}$ & $1.1\cdot10^{-5}$\\
$T_p$ (\PH) & $6.4\cdot10^{-13}$ & $1.2\cdot10^{-11}$ & $1.6\cdot10^{-11}$ & $1.5\cdot10^{-10}$ & $2.7\cdot10^{-10}$ & $6.7\cdot10^{-8}$\\
\bottomrule
\end{tabular}

\end{table}

\subsection{Thick cylinder: accuracy, convergence and stability}\label{sec:cylinder}
\paragraph{Polynomial field.} \Cref{tab:cylpoly} gives the maximum nodal error for the degree-5
field over all meshes and $p=5,\dots,8$. Both unsymmetric elements are exact in exact arithmetic;
in floating point the error of \PM\ grows with the distortion to $1.1\cdot10^{-4}$ ($45^\circ$) and
$1.6\cdot10^{-2}$ (folded mesh $90^\circ$), that of \PH\ stays $\le9.9\cdot10^{-10}$ on every mesh on which
it can be built. With exact integration (\PM: $n_q\ge p^2+p/2$ points) the residual of the exact
nodal values is $9.9\cdot10^{-7}$ for \PM\ and $5.0\cdot10^{-11}$ for \PH\ ($\alpha=60^\circ$, $N=2$,
$p=8$), with $\sigma_{\min}/\sigma_{\max}(\bm K)=1.1\cdot10^{-15}$, $4.8\cdot10^{-8}$ and $2.9\cdot10^{-6}$ for
\PM, \PH\ and \PP.

\begin{table}[t]
  \centering\small
  \caption{Thick cylinder, polynomial field (degree 5): maximum relative nodal error over $N=1,\dots,8$
    and $p=5,\dots,8$; in parentheses the number of rejected cases (singular interpolation matrix).}
  \label{tab:cylpoly}
  \begin{tabular}{@{}lrrrrrr@{}}
\toprule
 & $\alpha=0^\circ$ & $\alpha=15^\circ$ & $\alpha=30^\circ$ & $\alpha=45^\circ$ & $\alpha=60^\circ$ & $\alpha=90^\circ$\\
\midrule
\PM & $1.6\cdot10^{-11}$ & $3.8\cdot10^{-10}$ & $7.6\cdot10^{-9}$ & $1.1\cdot10^{-4}$ & $1.2\cdot10^{-5}$ (4) & $1.6\cdot10^{-2}$ (3)\\
\PH & $4.3\cdot10^{-13}$ & $3.3\cdot10^{-13}$ & $5.1\cdot10^{-13}$ & $9.9\cdot10^{-10}$ & $2.8\cdot10^{-10}$ (3) & $9.7\cdot10^{-10}$ (3)\\
\bottomrule
\end{tabular}

\end{table}

\paragraph{Convergence.} \Cref{fig:cyl-h} shows the energy error of the analytic field and of
the Lam\'e field at $\alpha=30^\circ$, \cref{tab:rates} the rates on the finest pair of meshes. On
these non-asymptotically-affine meshes the conventional element converges with
$\approx\lfloor p/2\rfloor$, because the parametric space contains $P_k(\bx)$ only for $2k\le p$
\cite{Arnold2002Approximation}. \PM\ is optimal for $\alpha=15^\circ$, $p\le6$, but at
$\alpha=30^\circ$ only for $p\le3$; beyond, its error is erratic and can increase under refinement
(analytic field, $p=4$: $\eE=2.2\cdot10^{-4}$ at $N=4$, $4.3\cdot10^{-4}$ at $N=8$). \PH\ converges with the
optimal order $p$ for all $p$ at $15^\circ$ and $30^\circ$ (Lam\'e: $1.96$--$7.87$). At $N=8$,
$\alpha=30^\circ$ (analytic field), its energy error is $8$--$2.6\cdot10^4$ times smaller than that of \PP,
and $1.4$--$6.7\cdot10^4$ times smaller than that of \PM\ for $p\ge3$ (equal at $p=2$). Under
$p$-refinement (\cref{fig:cyl-p}) \PH\ converges exponentially to $3.5\cdot10^{-13}$ ($30^\circ$) and
$1.1\cdot10^{-11}$ ($45^\circ$, folded elements), \PM\ only for $\alpha\le15^\circ$.

\begin{figure}[t]
  \centering
  \begin{tikzpicture}
    \begin{groupplot}[group style={group size=3 by 2, horizontal sep=3mm, vertical sep=9mm,
        yticklabels at=edge left, xticklabels at=edge bottom}, conv, width=5.0cm, height=4.6cm,
        ymin=1e-14, ymax=1, xlabel={$h\propto1/N$}]
      \nextgroupplot[title={\PM, analytic}, ylabel={$\eE$},
        legend to name=leg:p, legend columns=7, legend style={font=\footnotesize}]
      \addplot[pline=2] table[x=h, y=eE2]{\data/cyl_h_smooth_PM_a030.dat};
      \addplot[pline=3] table[x=h, y=eE3]{\data/cyl_h_smooth_PM_a030.dat};
      \addplot[pline=4] table[x=h, y=eE4]{\data/cyl_h_smooth_PM_a030.dat};
      \addplot[pline=5] table[x=h, y=eE5]{\data/cyl_h_smooth_PM_a030.dat};
      \addplot[pline=6] table[x=h, y=eE6]{\data/cyl_h_smooth_PM_a030.dat};
      \addplot[pline=7] table[x=h, y=eE7]{\data/cyl_h_smooth_PM_a030.dat};
      \addplot[pline=8] table[x=h, y=eE8]{\data/cyl_h_smooth_PM_a030.dat};
      \legend{$p=2$, $p=3$, $p=4$, $p=5$, $p=6$, $p=7$, $p=8$}
      \nextgroupplot[title={\PH, analytic}]
      \addplot[pline=2] table[x=h, y=eE2]{\data/cyl_h_smooth_PH_a030.dat};
      \addplot[pline=3] table[x=h, y=eE3]{\data/cyl_h_smooth_PH_a030.dat};
      \addplot[pline=4] table[x=h, y=eE4]{\data/cyl_h_smooth_PH_a030.dat};
      \addplot[pline=5] table[x=h, y=eE5]{\data/cyl_h_smooth_PH_a030.dat};
      \addplot[pline=6] table[x=h, y=eE6]{\data/cyl_h_smooth_PH_a030.dat};
      \addplot[pline=7] table[x=h, y=eE7]{\data/cyl_h_smooth_PH_a030.dat};
      \addplot[pline=8] table[x=h, y=eE8]{\data/cyl_h_smooth_PH_a030.dat};
      \nextgroupplot[title={\PP, analytic}]
      \addplot[pline=2] table[x=h, y=eE2]{\data/cyl_h_smooth_PP_a030.dat};
      \addplot[pline=3] table[x=h, y=eE3]{\data/cyl_h_smooth_PP_a030.dat};
      \addplot[pline=4] table[x=h, y=eE4]{\data/cyl_h_smooth_PP_a030.dat};
      \addplot[pline=5] table[x=h, y=eE5]{\data/cyl_h_smooth_PP_a030.dat};
      \addplot[pline=6] table[x=h, y=eE6]{\data/cyl_h_smooth_PP_a030.dat};
      \addplot[pline=7] table[x=h, y=eE7]{\data/cyl_h_smooth_PP_a030.dat};
      \addplot[pline=8] table[x=h, y=eE8]{\data/cyl_h_smooth_PP_a030.dat};
      \nextgroupplot[title={\PM, Lam\'e}, ylabel={$\eE$}]
      \addplot[pline=2] table[x=h, y=eE2]{\data/cyl_h_lame_PM_a030.dat};
      \addplot[pline=3] table[x=h, y=eE3]{\data/cyl_h_lame_PM_a030.dat};
      \addplot[pline=4] table[x=h, y=eE4]{\data/cyl_h_lame_PM_a030.dat};
      \addplot[pline=5] table[x=h, y=eE5]{\data/cyl_h_lame_PM_a030.dat};
      \addplot[pline=6] table[x=h, y=eE6]{\data/cyl_h_lame_PM_a030.dat};
      \addplot[pline=7] table[x=h, y=eE7]{\data/cyl_h_lame_PM_a030.dat};
      \addplot[pline=8] table[x=h, y=eE8]{\data/cyl_h_lame_PM_a030.dat};
      \nextgroupplot[title={\PH, Lam\'e}]
      \addplot[pline=2] table[x=h, y=eE2]{\data/cyl_h_lame_PH_a030.dat};
      \addplot[pline=3] table[x=h, y=eE3]{\data/cyl_h_lame_PH_a030.dat};
      \addplot[pline=4] table[x=h, y=eE4]{\data/cyl_h_lame_PH_a030.dat};
      \addplot[pline=5] table[x=h, y=eE5]{\data/cyl_h_lame_PH_a030.dat};
      \addplot[pline=6] table[x=h, y=eE6]{\data/cyl_h_lame_PH_a030.dat};
      \addplot[pline=7] table[x=h, y=eE7]{\data/cyl_h_lame_PH_a030.dat};
      \addplot[pline=8] table[x=h, y=eE8]{\data/cyl_h_lame_PH_a030.dat};
      \nextgroupplot[title={\PP, Lam\'e}]
      \addplot[pline=2] table[x=h, y=eE2]{\data/cyl_h_lame_PP_a030.dat};
      \addplot[pline=3] table[x=h, y=eE3]{\data/cyl_h_lame_PP_a030.dat};
      \addplot[pline=4] table[x=h, y=eE4]{\data/cyl_h_lame_PP_a030.dat};
      \addplot[pline=5] table[x=h, y=eE5]{\data/cyl_h_lame_PP_a030.dat};
      \addplot[pline=6] table[x=h, y=eE6]{\data/cyl_h_lame_PP_a030.dat};
      \addplot[pline=7] table[x=h, y=eE7]{\data/cyl_h_lame_PP_a030.dat};
      \addplot[pline=8] table[x=h, y=eE8]{\data/cyl_h_lame_PP_a030.dat};
    \end{groupplot}
  \end{tikzpicture}\\[1mm]
  \ref{leg:p}
  \caption{Thick cylinder, $\alpha=30^\circ$ (no folded elements): energy error versus element size
    for the analytic field (top) and the Lam\'e field (bottom).}
  \label{fig:cyl-h}
\end{figure}

\begin{table}[t]
  \centering\small
  \caption{Thick cylinder: observed rates of $\eE$ between $N=4$ and $N=8$ (optimal: $p$);
    ``--'': error at round-off level.}
  \label{tab:rates}
  \begin{tabular}{@{}llrrrrrrr@{}}
\toprule
 & & $p=2$ & $p=3$ & $p=4$ & $p=5$ & $p=6$ & $p=7$ & $p=8$\\
\midrule
analytic, $\alpha=15^\circ$ & \PM & 2.00 & 2.99 & 3.99 & 4.99 & 5.98 & -- & --\\
 & \PH & 2.00 & 3.00 & 4.00 & 4.99 & 6.01 & -- & --\\
 & \PP & 1.08 & 1.23 & 2.21 & 2.03 & 2.94 & 3.39 & 4.05\\
\addlinespace[2pt]
analytic, $\alpha=30^\circ$ & \PM & 2.00 & 2.86 & -0.98 & 2.40 & -- & -- & --\\
 & \PH & 2.00 & 3.00 & 4.01 & 5.00 & 6.02 & -- & --\\
 & \PP & 1.02 & 1.06 & 2.01 & 1.98 & 2.98 & 3.11 & 3.94\\
\addlinespace[2pt]
Lam\'e, $\alpha=15^\circ$ & \PM & 1.97 & 2.91 & 3.90 & 4.81 & 5.83 & -- & --\\
 & \PH & 1.96 & 2.94 & 3.94 & 4.92 & 5.91 & 6.88 & 7.87\\
 & \PP & 1.07 & 1.43 & 2.16 & 2.81 & 3.49 & 4.19 & 4.78\\
\addlinespace[2pt]
Lam\'e, $\alpha=30^\circ$ & \PM & 1.97 & 2.94 & 0.76 & 2.99 & 4.56 & -- & --\\
 & \PH & 1.96 & 2.96 & 3.94 & 4.94 & 5.92 & 6.90 & 7.87\\
 & \PP & 1.03 & 1.21 & 2.07 & 2.51 & 3.20 & 3.85 & 4.42\\
\bottomrule
\end{tabular}

\end{table}

\begin{figure}[t]
  \centering
  \begin{tikzpicture}
    \begin{groupplot}[group style={group size=3 by 1, horizontal sep=3mm, yticklabels at=edge left},
        width=5.0cm, height=4.6cm, ymode=log, xmin=2, xmax=8, ymin=1e-14, ymax=1e-1,
        xlabel={$p$}, unbounded coords=jump]
      \nextgroupplot[title={$\alpha=15^\circ$}, ylabel={$e_n$ ($N=8$)},
        legend style={font=\scriptsize, at={(0.03,0.03)}, anchor=south west}]
      \addplot[sPP] table[x=p, y=eNN8]{\data/cyl_p_smooth_PP_a015.dat};
      \addplot[sPM] table[x=p, y=eNN8]{\data/cyl_p_smooth_PM_a015.dat};
      \addplot[sPH] table[x=p, y=eNN8]{\data/cyl_p_smooth_PH_a015.dat};
      \legend{PP, PM, PH}
      \nextgroupplot[title={$\alpha=30^\circ$}]
      \addplot[sPP] table[x=p, y=eNN8]{\data/cyl_p_smooth_PP_a030.dat};
      \addplot[sPM] table[x=p, y=eNN8]{\data/cyl_p_smooth_PM_a030.dat};
      \addplot[sPH] table[x=p, y=eNN8]{\data/cyl_p_smooth_PH_a030.dat};
      \nextgroupplot[title={$\alpha=45^\circ$ (folded)}]
      \addplot[sPP] table[x=p, y=eNN8]{\data/cyl_p_smooth_PP_a045.dat};
      \addplot[sPM] table[x=p, y=eNN8]{\data/cyl_p_smooth_PM_a045.dat};
      \addplot[sPH] table[x=p, y=eNN8]{\data/cyl_p_smooth_PH_a045.dat};
    \end{groupplot}
  \end{tikzpicture}
  \caption{Thick cylinder, analytic field: $p$-convergence of the nodal error on the finest
    mesh.}
  \label{fig:cyl-p}
\end{figure}

\paragraph{Approximation versus stability.} Inserting the exact nodal values into the trial
functions gives the interpolation error of each space; the ratio of discretisation to
interpolation error isolates the stability of the discrete problem (\cref{tab:quasiopt}). For \PH\ it
lies between $0.78$ and $1.00$ in all $158$ cases on the meshes without folds with errors above the
round-off level, for \PP\ between $0.84$ and $1.00$ ($165$ cases): both are quasi-optimal with a
constant close to one. For \PM\ it ranges from $0.07$ to $10^5$; at $\alpha=30^\circ$, $p\ge3$ the
discrete \PM\ solution is up to $64$ times worse \emph{or} up to five times better than its
interpolant, i.e.\ it does not inherit the approximation property of its space. The \PM\ space
itself also approximates worse where its Lebesgue constant is large: at $\alpha=30^\circ$, $p\le6$, the
\PM\ interpolant of the analytic field is up to nine times less accurate than the \PH\
interpolant.

\begin{table}[t]
  \centering\small
  \caption{Thick cylinder, $N=8$: energy error of the discrete solution divided by the energy
    error of the interpolant of the exact field in the same trial space.}
  \label{tab:quasiopt}
  \begin{tabular}{@{}llrrrrrrr@{}}
\toprule
 & & $p=2$ & $p=3$ & $p=4$ & $p=5$ & $p=6$ & $p=7$ & $p=8$\\
\midrule
analytic, $\alpha=15^\circ$ & \PM & 0.99 & 0.98 & 1.00 & 0.98 & 1.01 & -- & --\\
 & \PH & 0.99 & 0.97 & 0.98 & 0.96 & 0.98 & -- & --\\
 & \PP & 0.97 & 0.94 & 0.96 & 0.93 & 0.96 & 0.95 & 0.96\\
\addlinespace[2pt]
analytic, $\alpha=30^\circ$ & \PM & 0.91 & 0.20 & 64.38 & 3.40 & 0.38 & -- & --\\
 & \PH & 0.96 & 0.93 & 0.93 & 0.89 & 0.90 & -- & --\\
 & \PP & 0.94 & 0.92 & 0.92 & 0.89 & 0.93 & 0.93 & 0.93\\
\addlinespace[2pt]
Lam\'e, $\alpha=15^\circ$ & \PM & 1.00 & 0.99 & 0.99 & 0.98 & 0.99 & -- & --\\
 & \PH & 0.98 & 0.95 & 0.96 & 0.95 & 0.96 & 0.94 & 0.95\\
 & \PP & 0.98 & 0.96 & 0.98 & 0.97 & 0.98 & 0.97 & 0.98\\
\addlinespace[2pt]
Lam\'e, $\alpha=30^\circ$ & \PM & 0.99 & 0.48 & 14.98 & 5.46 & 8.75 & -- & --\\
 & \PH & 0.93 & 0.91 & 0.92 & 0.89 & 0.90 & 0.88 & 0.89\\
 & \PP & 0.95 & 0.93 & 0.96 & 0.93 & 0.95 & 0.94 & 0.95\\
\bottomrule
\end{tabular}

\end{table}

\paragraph{Lebesgue constants, conditioning and spectra.} \Cref{tab:lebesgue} lists the
maxima over the elements of the finest mesh. The Lebesgue constant of \PM\ reaches $10^3$ at
$30^\circ$, $p=4$ and $1.3\cdot10^8$ at $45^\circ$, $p=8$; that of \PH\ stays within $30\,\%$ of the GLL value
up to $15^\circ$ and $\le95$ at $30^\circ$. $\kappa_2(\bm V)$ of \PM\ exceeds its Lebesgue constant by up
to two orders of magnitude (extrapolation at the reference points, \cref{prop:kappa}); that of \PH\
is $10^2$--$10^5$ times smaller than that of \PM\ at $30^\circ$, $p\ge6$. At mesh level, the smallest singular values of the global
stiffness matrices of the three elements agree to $1\,\%$ (14 cases, $\alpha\le30^\circ$), but the
largest singular value of \PM\ exceeds that of \PP\ by up to a factor $344$ ($30^\circ$, $N=4$,
$p=6$), whereas that of \PH\ stays within $1.2$ of it. In the expansion of the error of the analytic
field in the singular triplets of $\bm K$, a single one of the six lowest triplets carries
$21$--$68\,\%$ of the \PM\ error in four of five cases with $\alpha\ge30^\circ$; for \PH\ and \PP\ the error
is spread over the spectrum (largest single share $2$--$8\,\%$), and the lowest $50$ singular values
of \PH\ agree with those of \PP\ to $2.3\,\%$.

\begin{table}[t]
  \centering\small
  \caption{Thick cylinder, $N=8$: Lebesgue constant $\Lambda$ ($16\times16$ sample points per
    element, a lower bound) and condition number $\kappa_2(\bm V)$ of the interpolation matrix,
    maximum over the elements.}
  \label{tab:lebesgue}
  \begin{tabular}{@{}llrrrrrrr@{}}
\toprule
 & & $p=2$ & $p=3$ & $p=4$ & $p=5$ & $p=6$ & $p=7$ & $p=8$\\
\midrule
$\Lambda$, $\alpha=15^\circ$ & \PM & $2.5$ & $6.0$ & $1.5\cdot10^{1}$ & $3.7\cdot10^{1}$ & $1.3\cdot10^{2}$ & $1.2\cdot10^{4}$ & $1.3\cdot10^{4}$\\
 & \PH & $1.8$ & $2.4$ & $2.8$ & $3.1$ & $3.5$ & $4.5$ & $4.2$\\
$\Lambda$, $\alpha=30^\circ$ & \PM & $1.0\cdot10^{1}$ & $2.4\cdot10^{2}$ & $1.1\cdot10^{3}$ & $1.6\cdot10^{3}$ & $6.2\cdot10^{3}$ & $4.2\cdot10^{6}$ & $1.2\cdot10^{6}$\\
 & \PH & $2.0$ & $4.8$ & $7.5$ & $1.6\cdot10^{1}$ & $2.5\cdot10^{1}$ & $7.4\cdot10^{1}$ & $9.5\cdot10^{1}$\\
$\Lambda$, $\alpha=45^\circ$ & \PM & $1.9\cdot10^{2}$ & $1.5\cdot10^{3}$ & $1.3\cdot10^{4}$ & $5.0\cdot10^{5}$ & $8.3\cdot10^{5}$ & $1.4\cdot10^{6}$ & $1.3\cdot10^{8}$\\
 & \PH & $2.8$ & $1.0\cdot10^{1}$ & $2.1\cdot10^{1}$ & $4.7\cdot10^{1}$ & $9.1\cdot10^{1}$ & $2.6\cdot10^{2}$ & $5.5\cdot10^{2}$\\
 & GLL & $1.56$ & $2.22$ & $2.68$ & $3.00$ & $3.41$ & $3.47$ & $4.15$\\
\midrule
$\kappa_2(\bm V)$, $\alpha=15^\circ$ & \PM & $4.8$ & $1.2\cdot10^{1}$ & $6.0\cdot10^{1}$ & $2.9\cdot10^{2}$ & $1.9\cdot10^{3}$ & $3.1\cdot10^{5}$ & $2.4\cdot10^{5}$\\
 & \PH & $5.5$ & $9.7$ & $1.9\cdot10^{1}$ & $4.0\cdot10^{1}$ & $9.4\cdot10^{1}$ & $2.3\cdot10^{2}$ & $6.0\cdot10^{2}$\\
$\kappa_2(\bm V)$, $\alpha=30^\circ$ & \PM & $3.8\cdot10^{1}$ & $1.7\cdot10^{3}$ & $5.9\cdot10^{3}$ & $8.5\cdot10^{3}$ & $8.6\cdot10^{4}$ & $2.6\cdot10^{8}$ & $2.6\cdot10^{8}$\\
 & \PH & $9.0$ & $2.4\cdot10^{1}$ & $6.0\cdot10^{1}$ & $1.8\cdot10^{2}$ & $7.0\cdot10^{2}$ & $2.7\cdot10^{3}$ & $9.8\cdot10^{3}$\\
$\kappa_2(\bm V)$, $\alpha=45^\circ$ & \PM & $8.4\cdot10^{2}$ & $1.5\cdot10^{4}$ & $2.1\cdot10^{5}$ & $2.9\cdot10^{7}$ & $2.8\cdot10^{8}$ & $4.7\cdot10^{9}$ & $7.9\cdot10^{11}$\\
 & \PH & $1.8\cdot10^{1}$ & $9.1\cdot10^{1}$ & $4.1\cdot10^{2}$ & $2.4\cdot10^{3}$ & $3.3\cdot10^{4}$ & $3.2\cdot10^{5}$ & $4.0\cdot10^{6}$\\
\bottomrule
\end{tabular}

\end{table}

\paragraph{Limits of the hybrid space.} (i) On the folded meshes with $\alpha\ge60^\circ$ the
interpolation matrix of \PH\ becomes singular in $6$ of the $24$ cases with $p\ge6$ (that of \PM\ in
$7$). \Cref{tab:unisolvence} confirms \cref{prop:factor}: $\mathrm{rcond}(\bm A)$ follows
$\mathrm{rcond}(\bm V_T)$ within a factor $30$, the determinant identity~\eqref{eq:factor} holds to
$\le3\cdot10^{-13}$ (relative) for $\mathrm{rcond}(\bm V_T)\ge5\cdot10^{-6}$ and degrades with the
conditioning as expected for computed determinants, and the singular cases are exactly those with
$\mathrm{rcond}(\bm A)\le4\cdot10^{-16}$ (all other cases: $\ge2\cdot10^{-15}$). (ii) On the same meshes the \PH\ error for $p=2$ ($60^\circ$)
and $p\le4$ ($90^\circ$) is non-monotone in $N$ and of order $10^{-2}$--$1$, as for the other two
elements; the singular complement integrand (\cref{sec:hybrid}) is one possible cause. (iii) The
error constant of \PH\ depends on the distortion: for the Lam\'e field at $\alpha=15^\circ$, where \PM\ is
stable, \PH\ is up to $4.1$ times less accurate than \PM\ ($p\le7$), because its interpolation error
is $1.2$--$5.7$ times larger ($N=8$); the parametric complement approximates the $r^{-2}$ terms of this
field worse than the full $Q_p$ in skew coordinates.

\begin{table}[t]
  \centering\small
  \caption{Unisolvence of the hybrid space on the thick-cylinder meshes: minimum over the elements
    of $\mathrm{rcond}$ of $\bm V_T$, of $\bm A$ (\cref{prop:factor}) and of $\bm V$ of
    $Q_p(\hat x,\hat y)$; number of elements with $\mathrm{rcond}(\bm V_T)<\varepsilon$.}
  \label{tab:unisolvence}
  \begin{tabular}{@{}rrrcccc@{}}
\toprule
$\alpha$ & $N$ & $p$ & $\min\mathrm{rcond}(\bm V_T)$ & $\min\mathrm{rcond}(\bm A)$ & $\min\mathrm{rcond}(\bm V_Q)$ & singular\\
\midrule
$0^\circ$ & 8 & 8 & $1.9\cdot10^{-3}$ & $5.5\cdot10^{-2}$ & $1.0\cdot10^{-2}$ & 0\\
$15^\circ$ & 8 & 8 & $7.5\cdot10^{-5}$ & $5.0\cdot10^{-4}$ & $4.0\cdot10^{-7}$ & 0\\
$30^\circ$ & 8 & 6 & $1.0\cdot10^{-4}$ & $5.5\cdot10^{-4}$ & $2.9\cdot10^{-6}$ & 0\\
$30^\circ$ & 8 & 7 & $2.1\cdot10^{-5}$ & $1.1\cdot10^{-4}$ & $6.3\cdot10^{-10}$ & 0\\
$30^\circ$ & 8 & 8 & $5.3\cdot10^{-6}$ & $2.7\cdot10^{-5}$ & $6.3\cdot10^{-10}$ & 0\\
$45^\circ$ & 8 & 6 & $2.8\cdot10^{-6}$ & $1.3\cdot10^{-5}$ & $1.2\cdot10^{-9}$ & 0\\
$45^\circ$ & 8 & 8 & $1.8\cdot10^{-8}$ & $8.1\cdot10^{-8}$ & $4.7\cdot10^{-13}$ & 0\\
$60^\circ$ & 4 & 8 & $1.4\cdot10^{-17}$ & $9.8\cdot10^{-17}$ & $7.5\cdot10^{-19}$ & 4\\
$60^\circ$ & 8 & 6 & $3.1\cdot10^{-16}$ & $2.2\cdot10^{-15}$ & $1.3\cdot10^{-16}$ & 0\\
$60^\circ$ & 8 & 7 & $5.2\cdot10^{-19}$ & $6.5\cdot10^{-18}$ & $1.7\cdot10^{-18}$ & 8\\
$60^\circ$ & 8 & 8 & $3.4\cdot10^{-20}$ & $3.1\cdot10^{-19}$ & $6.1\cdot10^{-21}$ & 8\\
$90^\circ$ & 1 & 7 & $7.8\cdot10^{-16}$ & $7.8\cdot10^{-15}$ & $1.5\cdot10^{-15}$ & 0\\
$90^\circ$ & 1 & 8 & $6.1\cdot10^{-18}$ & $6.1\cdot10^{-17}$ & $5.3\cdot10^{-19}$ & 1\\
$90^\circ$ & 8 & 6 & $1.2\cdot10^{-14}$ & $9.3\cdot10^{-14}$ & $6.0\cdot10^{-15}$ & 0\\
$90^\circ$ & 8 & 7 & $3.9\cdot10^{-17}$ & $3.9\cdot10^{-16}$ & $2.8\cdot10^{-17}$ & 8\\
$90^\circ$ & 8 & 8 & $3.9\cdot10^{-20}$ & $9.4\cdot10^{-19}$ & $1.2\cdot10^{-19}$ & 16\\
\bottomrule
\end{tabular}

\end{table}
